
\documentclass[runningheads]{llncs}
\usepackage{graphicx}
\usepackage[utf8]{inputenc}
\usepackage{fancyhdr}
 
\pagestyle{fancy}
\fancyhf{}
\fancyhf{}
\rhead{ESPE}
\lhead{Universidad de las Fuerzas Armadas}
\rfoot{\thepage}

% If you use the hyperref package, please uncomment the following line
% to display URLs in blue roman font according to Springer's eBook style:
% \renewcommand\UrlFont{\color{blue}\rmfamily}

\begin{document}
%
\title{Titulo del Proyecto}

%\titlerunning{Abbreviated paper title}
% If the paper title is too long for the running head, you can set
% an abbreviated paper title here
%
\author{Integrante 1\inst{1}\orcidID{L00374562} \and
Integrante 2\inst{2}\orcidID{L00741236} \and
Integrante 3\inst{3}\orcidID{L00123456}}
%
\authorrunning{Universidad de las Fuerzas Armadas}
% First names are abbreviated in the running head.
% If there are more than two authors, 'et al.' is used.
%
\institute{Universidad de las Fuerzas Armadas\\
\inst{1}\email{integrante1@espe.edu.ec}\\
\inst{2}\email{integrante2@espe.edu.ec}\\
\inst{3}\email{integrante3@espe.edu.ec}}
%
\maketitle 
%%%%%%%%%%%%%%%%%%%%%%%%%%%%%%%%%%%%%%%%%%%%%%%%%%%%%%%%%%%%%%%%%
%%                            ABSTRACT                         %%
%%%%%%%%%%%%%%%%%%%%%%%%%%%%%%%%%%%%%%%%%%%%%%%%%%%%%%%%%%%%%%%%%
\begin{abstract}
The abstract should briefly summarize the contents of the paper in
15--250 words.
\keywords{First keyword  \and Second keyword \and Another keyword.}
\end{abstract}
%
%%%%%%%%%%%%%%%%%%%%%%%%%%%%%%%%%%%%%%%%%%%%%%%%%%%%%%%%%%%%%%%%%
%%                        INTRODUCCIÓN                         %%
%%%%%%%%%%%%%%%%%%%%%%%%%%%%%%%%%%%%%%%%%%%%%%%%%%%%%%%%%%%%%%%%%
\section{Introducción}
La introducción es una descripción sencilla, directa y concisa de las ideas principales del proyecto. El objetivo de una introducción de un proyecto será identificar y transmitir el tema y los resultados que esperamos alcanzar al término del proyecto. Entre otros detalles como la manera en la que trabajaremos y las herramientas que utilizaremos.

Para poder llegar a cumplir con estos objetivos, es importante presentar esta introducción de la manera más simple y directa posible. Recogemos así las características que no deben faltar en una adecuada introducción para un proyecto.
\begin{itemize}
    \item Identificación del objetivo principal y la calidad de los resultados esperados, además estos deben ser escritos en prosa.
    \item Comunicación clara y ordenada de los procedimientos a seguir.
    \item Una introducción de un proyecto plantea reflexiones y dudas que se irán resolviendo a lo largo del mismo.
\end{itemize}
Como se menciono antes, la introducción reúne gran parte de información de todo el proyecto, por lo que se debe indicar en un párrafo cual es la estructura del artículo indicando que se trata en cada sección.
%%%%%%%%%%%%%%%%%%%%%%%%%%%%%%%%%%%%%%%%%%%%%%%%%%%%%%%%%%%%%%%%%
%%                      ESTADO DEL ARTE                        %%
%%%%%%%%%%%%%%%%%%%%%%%%%%%%%%%%%%%%%%%%%%%%%%%%%%%%%%%%%%%%%%%%%
\section{Estado del Arte}
El estado del arte describe las investigaciones más recientes y actuales que sobre un tema en específico se han
realizado.
La descripción es un texto académico que expone sistemáticamente los avances existentes acerca de un tema y es de carácter más cualitativo, en el que se detallan los resultados y enfoques de las investigaciones en torno al tema que cada investigación ha abonado al tema de estudio de interés del investigador que elabora el estado del arte. La descripción gira en torno esencial a cuatro elementos:
\begin{itemize}
    \item Quien? El investigador que desarrolló estudio.
    \item Cuándo? El año en que se publicaron los resultados del estudio. Aunque sabemos anticipadamente que el estudio debió de haber sido desarrollado con anterioridad mínima a un año generalmente.
    \item ¿Qué? El objeto de estudio. Es aquí en donde se hace énfasis en la descripción. No solo se dice el objeto de estudio, sino el enfoque, los resultados de la investigación.
    \item ¿Dónde? El lugar en donde se realizó la investigación. Este es un dato de referencia con varios propósitos: uno es para organizar la información de lo macro a micro (de carácter internacional, nacional o local); otro propósito es para saber la manera de establecer contacto con el autor de la investigación si así fuera el deseo del investigador que realiza el estado del arte; por ejemplo si es de la localidad puede contactarlo de manera directa y cara a cara, si no tendrá que establecer contacto por otros medios, ahora tenemos al alcance los medios electrónicos para ello que recortan el tiempo de la retroalimentación de un mensaje.
\end{itemize}
%%%%%%%%%%%%%%%%%%%%%%%%%%%%%%%%%%%%%%%%%%%%%%%%%%%%%%%%%%%%%%%%%
%%                        MARCO TEÓRICO                        %%
%%%%%%%%%%%%%%%%%%%%%%%%%%%%%%%%%%%%%%%%%%%%%%%%%%%%%%%%%%%%%%%%%
\section{Marco Teórico}
Este punto requiere que el estudiante realice una amplia consulta bibliográfica sobre el tema de su trabajo.
Se describe la teoría o conjunto teórico apropiado con la cual el investigador enfrenta su proyecto y la realidad dentro del cual se ubica el problema de investigación, incluye:
\begin{itemize}
    \item Describir la actual relación entre el problema enunciado y el sistema o sistemas teóricos/conceptuales que pueden guiarlo.
    \item Dejar claramente especificada la relación entre la teoría que guía la investigación y la realidad que se percibe como problema de investigación.
    \item Conceptualizar el problema en la forma de un modelo, útil para clarificar los conceptos y relaciones conceptuales.
    \item Señalar la forma en que la investigación actual enriquece, amplía y profundiza el conocimiento teórico, sustantivo y metodológico acumulado en estudios previos. Para su desarrollo debe evitarse en lo posible redactar párrafos continuos de texto, para lo cual se debe utilizar mentefactos, cuadro comparativos, esquemas o mapas conceptuales
\end{itemize}
\textbf{Nota: }El marco teórico no puede exceder más de una página.
%%%%%%%%%%%%%%%%%%%%%%%%%%%%%%%%%%%%%%%%%%%%%%%%%%%%%%%%%%%%%%%%%
%%                  EXPLICACIÓN DEL CÓDIGO                     %%
%%%%%%%%%%%%%%%%%%%%%%%%%%%%%%%%%%%%%%%%%%%%%%%%%%%%%%%%%%%%%%%%%
\section{Diseño}
En esta sección se explicara el diseño que se planteo antes de implementar el proyecto. 
\section{Implementación}
Describir brevemente como se desarrollo el proyecto, tanto en hardware como en software.
\section{Pruebas de Funcionamiento}
Presentar las pruebas que apoyan tales resultados cumplidos de los objetivos, sea en forma de figuras, tablas o en el mismo texto. 
%%%%%%%%%%%%%%%%%%%%%%%%%%%%%%%%%%%%%%%%%%%%%%%%%%%%%%%%%%%%%%%%%
%%                       CONCLUSIONES                          %%
%%%%%%%%%%%%%%%%%%%%%%%%%%%%%%%%%%%%%%%%%%%%%%%%%%%%%%%%%%%%%%%%%
\section{Conclusiones}
Se establece las conclusiones de cada asunto investigado, implicaciones para la teoría y resultados de las
experiencias. Estos siempre estarán en relaciona los objetivos generales y específicos.

%%%%%%%%%%%%%%%%%%%%%%%%%%%%%%%%%%%%%%%%%%%%%%%%%%%%%%%%%%%%%%%%%
%%                       BIBLIOGRAFÍA                          %%
%%%%%%%%%%%%%%%%%%%%%%%%%%%%%%%%%%%%%%%%%%%%%%%%%%%%%%%%%%%%%%%%%
% \bibliographystyle{splncs04}
% \bibliography{mybibliography}

\begin{thebibliography}{8}
\bibitem{ref_article1}
Author, F.: Article title. Journal \textbf{2}(5), 99--110 (2016)

\bibitem{ref_lncs1}
Author, F., Author, S.: Title of a proceedings paper. In: Editor,
F., Editor, S. (eds.) CONFERENCE 2016, LNCS, vol. 9999, pp. 1--13.
Springer, Heidelberg (2016). \doi{10.10007/1234567890}

\bibitem{ref_book1}
Author, F., Author, S., Author, T.: Book title. 2nd edn. Publisher,
Location (1999)

\bibitem{ref_proc1}
Author, A.-B.: Contribution title. In: 9th International Proceedings
on Proceedings, pp. 1--2. Publisher, Location (2010)

\bibitem{ref_url1}
LNCS Homepage, \url{http://www.springer.com/lncs}. Last accessed 4
Oct 2017
\end{thebibliography}
\end{document}
